\SourceLeaf{342}%
nimus proprie accepta), habebuntur aequationes, ex concl. 3 art. 235 facile deducendae:
\[
\begin{gathered}
p''hn'-p'h'n-p(in'-i'n)=0\\
(p''-p')(in'+i'n)-p(kn'-k'n)+p'''(hn'-h'n)=0\\
p'kn'-p''k'n-p'''(in'-i'n)=0
\end{gathered}
\]
tresque aliae ex his per commutationem numerorum $p$, $p'$, $p''$, $p'''$ cum $q$, $q'$, $q''$, $q'''$ oriundae; $n$, $n'$ sunt radices quadratae positivae e quotientibus prodeuntibus, si determinantes formarum $(h,i,k)$, $(h',i',k')$ per det. formae $F$ dividuntur. Si itaque hae formae sunt identicae, sive $n=n'$, $h=h'$, $i=i'$, $k=k'$, illae aequationes transeunt in has:
\[
(p''-p')hn=0,\quad (p''-p')in=0,\quad (p''-p')kn=0
\]
unde erit \textit{necessario} $p'=p''$, prorsusque simili modo $q'=q''$. --- Tribuendo itaque formis $(h,i,k)$, $(h',i',k')$ \textit{easdem} indeterminatas $t$, $u$, designandoque indeterminatas formae $F$ per $T$, $U$, transibit $F$ per substitutionem
\[
T=ptt+2p'tu+p'''uu,\quad U=qtt+2q'tu+q'''uu\quad\text{in}\quad (htt+2itu+kuu)^2
\]

II. Si forma $F$ oritur e duplicatione formae $f$, orietur etiam e duplicatione cuiusvis alius formae cum $f$ in eadem classe contentae sive classis formae $F$ e duplicatione classis formae $f$ (V. art. 238). Ita in ex. art. praec. $(5,2,31)$ orietur etiam e duplicatione formae $(11,-5,16)$, ipsi $(11,17,40)$ proprie aequivalentis. Ex una classe, per cuius dupl. classis formae $F$ oritur, \textit{omnes} (si plures dantur) inveniuntur adiumento probl. 260; in exemplo nostro alia huiusmodi classis positiva non dabitur, quia una tantummodo classis anceps proprie primitiva positiva det. $-151$ exstat (puta principalis); quum e compositione classis unicae ancipitis negativae $(-1,0,-151)$, cum classe $(11,-5,16)$ oriatur classis $(-11,-5.\ -16)$, haec erit unica negativa, e cuius duplicatione classis $(5,2,31)$ oritur.

III. Quum per solutionem ipsam probl. art. praec. evictum sit, quamvis classem formarum binariarum proprie primitivam (positivam) ad genus principale pertinentem ex alicuius classis pr. prim. eiusdem det. duplicatione oriri posse: theorema art. 261, per quod certi eramus, \textit{ad minimum} semissi omnium characterum
\SourceLeaf{343}
pro determinante non-quadrato dato $D$ assignabilium genera proprie primitiva (positiva) respondere non posse, eo iam ampliatur, ut \textit{praecise} semissi omnium horum characterum talia genera revera respondeant, alterique ideo semissi nulla respondere possint (V. demonstr. illius theor.). Quare quum in art. 264 omnes illi characteres assignabiles in duas species $P$, $Q$ aequaliter distributi sint, e quibus posteriores $Q$ formis pr. prim. (positivis) respondere non posse probatum erat, de reliquis autem $P$ incertum maneret, an singulis genera semper revera responderent: nunc hoc dubium penitus est sublatum, certique sumus, in toto characterum complexu $P$ nullum adesse, cui genus non respondeat. --- Hinc facile quoque deducitur, pro determinante negativo in ordine pr. prim. \textit{negativo}, in quo omnes $P$ impossibiles \textit{solosque} $Q$ possibiles esse in art. 264, I ostensum est. \textit{omnes} $Q$ revera possibiles esse. Designante enim $K$ characterem quemcunque ex $Q$, $f$ formam arbitrariam ex ordine pr. prim. neg. formarum det. $D$, atque $K'$ ipsius characterem, hic erit ex $Q$; unde facile perspicitur, characterem ex $K$, $K'$ compositum (ad normam art. 246) ad $P$ pertinere, adeoque formas pr. primitivas positivas det. $D$ exstare, quae ei respondeant; ex compositione talis formae cum $f$ manifesto orietur forma pr. prim. neg. det. $D$, cuius character erit $K$. --- Prorsus simili ratione probatur, in ordine improprie primitivo eos characteres, qui per praecepta art. 264 II, III \textit{soli} possibiles inveniuntur, \textit{omnes} possibiles esse, sive sint $P$ sive $Q$. --- Haecce theoremata, ni vehementer fallimur, ad pulcherrima in theoria formarum binariarum sunt referenda, eo magis quod licet summa simplicitate gaudeant, tamen tam recondita sint ut ipsarum demonstrationem rigorosam absque tot aliarum disquisitionum subsidio condere non liceat.

\Needspace{8\baselineskip}
\begin{center}
\textit{Theoria decompositionis tum numerorum tum formarum binariarum in tria quadrata.}
\end{center}
Transimus iam ad aliam applicationem digressionis praecedentis, ad discerptionem tum numerorum tum formarum binariarum in terna quadrata, cui praemittimus sequens

\Needspace{8\baselineskip}
\Article{288}
\textsc{Problema.} \textit{Designante $M$ numerum positivum, invenire conditiones sub quibus formae binariae primitivae negativae determinantis $-M$ dari possint, quae sint residua quadratica ipsius $M$ sive pro quibus $1$ sit numerus characteristicus.}

\SourceLeaf{344}
\textit{Sol.} Designemus per $\Omega$ complexum omnium characterum particularium, quos praebent relationes numeri $1$ tum ad singulos divisores primos (impares) ipsius $M$ tum ad numerum $8$ vel $4$, quando ipsum $M$ metitur; manifesto hi characteres erunt $Rp$, $Rp'$, $Rp''$ etc., denotantibus $p$, $p'$, $p''$ etc. illos divisores primos, atque $1,4$ quando $4$; $1,8$ quando $8$ ipsum $M$ metitur. Praeterea utamur literis $P$, $Q$ in eadem significatione ut in art. praec. sive ut in 264. Iam distinguamus casus sequentes.

I. Quando $M$ per $4$ divisibilis est, $\Omega$ erit character integer, patetque ex art. 233 V, $1$ talium tantummodo formarum numerum characteristicum esse posse, quarum character sit $\Omega$. Sed manifestum est, $\Omega$ fore characterem formae principalis $(1,0,M)$, adeoque ad $P$ pertinere et proin formae proprie primitivae negativae competere non posse; quare quum formae improprie primitivae pro tali det. non dentur, nullae omnino formae prim. neg. in hoc casu dantur, quae sint residua ipsius $M$.

II. Quando $M\equiv3\ (\mathrm{mod.}\ 4)$, prorsus eadem ratiocinia valent ea sola exceptione ut in hoc casu ordo \textit{improprie} primitivus negativus exstet, in quo characteres $P$ vel possibiles erunt, vel impossibiles, prout $M\equiv3$ vel $\equiv7\ (\mathrm{mod.}\ 8)$, V. art. 264, III. In casu igitur priori in hoc ordine genus dabitur, cuius character sit $\Omega$, unde $1$ erit numerus characteristicus omnium formarum in ipso contentarum; in casu posteriori nullae omnino formae negativae hac proprietate praeditae dari poterunt.

III. Quando $M\equiv1\ (\mathrm{mod.}\ 4)$, $\Omega$ nondum est character completus, sed insuper accedere debet relatio ad numerum $4$; patet autem, $\Omega$ necessario in characterem formae, cuius num. char. sit $1$, ingredi debere, et vice versa, formam quamvis, cuius character sit vel $\Omega;\ 1,4$, vel $\Omega;\ 3,4$, habere numerum char. $1$. Iam $\Omega;\ 1,4$ manifesto est character generis principalis, qui ad $P$ pertinet adeoque in ordine pr. prim. negativo impossibilis est; ex eadem ratione $\Omega;\ 3,4$ ad $Q$ pertinebit (art. 263), unde ipsi in ordine pr. prim. negativo genus respondebit, cuius formae omnes habebunt num. char. $1$. Ordo improprie primitivus in hoc casu, ut in sequente, non datur.

IV. Quando $M\equiv2\ (\mathrm{mod.}\ 4)$, ad $\Omega$ accedere debet relatio ad $8$, quo fiat character completus, puta vel $1\mathbin{\textit{et}}3,8$, vel $5\mathbin{\textit{et}}7,8$, quando $M\equiv2\ (\mathrm{mod.}\ 8)$; et vel $1\mathbin{\textit{et}}7,8$, vel $3\mathbin{\textit{et}}5,8$, quando $M\equiv6\ (\mathrm{mod.}\ 8)$. Pro casu priori character $\Omega;\ 1\mathbin{\textit{et}}3,\allowbreak 8$ manifesto pertinet ad $P$, adeoque $\Omega;\ 5\mathbin{\textit{et}}7,\allowbreak 8$ ad $Q$, unde ipsi respon%
\SourceLeaf{345}%
debit genus pr. prim. neg.; similique ratione pro posteriori unum genus in ordine pr. prim. negativo dabitur, cuius formae proprietate praescripta praeditae sint, puta cuius character $\Omega;\ 3\mathbin{\textit{et}}5,\allowbreak 8$.

Ex his colligitur, formas primitivas negativas det. $-M$, quarum numerus characteristicus sit $1$, dari, quando $M$ alicui numerorum $1$, $2$, $3$, $5$, $6$ secundum modulum $8$ congruus sit et quidem in unico semper genere, quod improprium erit quando $M\equiv3$; tales formas omnino non dari, quando $M\equiv0,\allowbreak \ 4\text{ vel }7\ (\mathrm{mod.}\ 8)$. Ceterum manifestum est, si $(-a,-b,-c)$ sit forma primitiva negativa, cuius num. char. $+1$, $(a,b,c)$ esse formam primitivam positivam, cuius num. char. $-1$; hinc perspicuum est, in quinque casibus prioribus (quando $M\equiv1,2,3,5,6$) dari genus unum primitivum positivum, cuius formae habeant num. char. $-1$, et quidem pro $M\equiv3$ \textit{improprium}, in tribus reliquis vero (quando $M\equiv0,4,7$) tales formas positivas omnino dari non posse.

\Needspace{8\baselineskip}
\Article{289}
Circa repraesentationes proprias formarum binariarum per ternariam $xx+yy+zz=f$, e theoria generali in art. 282 tradita colliguntur haec:

I. Forma binaria $\varphi$ per $f$ proprie repraesentari nequit, nisi fuerit forma positiva primitiva, atque $-1$ (\textit{i. e.} det. formae $f$) ipsius numerus characteristicus. Quare pro determinante positivo, nec non pro negativo $-M$, quando $M$ est vel per $4$ divisibilis vel formae $8n+7$, nullae formae binariae per $f$ proprie repraesentabiles dantur.

II. Si vero $\varphi=(p,q,r)$ est forma positiva primitiva determinantis $-M$, atque $-1$ numerus characteristicus formae $\varphi$, adeoque etiam oppositae $(p,-q,r)$: dabuntur repraesentationes propriae formae $\varphi$ per $f$ ad quemlibet valorem datum expr. $\surd{}-(p,-q,r)$ pertinentes. Scilicet omnes coëfficientes formae ternariae $g$ det. $-1$ (art. 283) necessario fient integri, $g$ vero forma definita, adeoque ipsi $f$ certo aequivalens (art. 285, I).

III. Multitudo omnium repraesentationum ad eundem valorem expr. $\surd{}-(p,-q,r)$ pertinentium in omnibus casibus, praeter $M=1$ et $M=2$, per art. 283, III aeque magna est ac multitudo transformationum formae $f$ in $g$, adeoque, per art. 285, $=48$; ibinde patet, si una repraesentatio ad valorem datum pertinens habeatur, $47$ reliquas inde derivari, valores ipsorum $x$, $y$, $z$
\SourceLeaf{346}
omnibus quibus fieri potest modis tum inter se permutando tum signis oppositis afficiendo; quare omnes $48$ repraesentationes \textit{unicam} decompositionem formae $\varphi$ in tria quadrata producunt, si ad quadrata ipsa tantum, neque ad ipsorum ordinem radicumve signa respicitur.

IV. Posita multitudine omnium numerorum primorum imparium diversorum ipsum $M$ metientium $=\mu$, haud difficile ex art. 233 concluditur, multitudinem omnium valorum diversorum expressionis $\surd{}-(p,\allowbreak -q,\allowbreak r)(\mathrm{mod.}\ M)$ fore $=2^\mu$, e quibus per art. 283 semissem tantum considerare oportet (quando $M>2$). Quare multitudo omnium repraesentationum propriarum formae $\varphi$ per $f$ erit $=48\,\allowbreak  .\,\allowbreak  2^{\mu-1}=3\,\allowbreak  .\,\allowbreak  2^{\mu+3}$; multitudo autem discerptionum diversarum in terna quadrata $=2^{\mu-1}$.

\textit{Ex.} Sit $\varphi=19tt+6tu+41uu$, adeoque $M=770$; hic quatuor valores sequentes expr. $\surd{}-(19,\allowbreak -3,\allowbreak 41)(\mathrm{mod.}\ 770)$ considerare oportet (art. 283): $(39,237)$, $(171,-27)$, $(269,-83)$, $(291,-127)$. Ut inveniantur repraesentationes ad valorem $(39,237)$ pertinentes, primo eruitur forma ternaria $\left(\begin{smallmatrix}19,&41,&2\\3,&6,&3\end{smallmatrix}\right)=g$, in quam per praecepta artt. 272, 275 $f$ transire invenitur per substitutionem $\left\{\begin{smallmatrix}1,&-6,&-0\\-3,&-2,&-1\\-3,&-1,&-1\end{smallmatrix}\right\}$ unde habetur repraesentatio formae $\varphi$ per $f$ haec:
\[
x=t-6u,\quad y=-3t-2u.\quad z=-3t-u
\]
repraesentationes $47$ reliquas ad eundem valorem pertinentes, quae ex horum valorum permutatione signorumque conversione oriuntur, brevitatis caussa non adscribimus. Omnes vero $48$ repraesentationes eandem discerptionem formae $\varphi$ in tria quadrata
\[
tt-12tu+36uu,\quad 9tt+12tu+4uu,\quad 9tt+6tu+uu
\]
producunt.

Prorsus simili modo valor $(171,-27)$ suppeditat discerptionem in quadrata $(3t+5u)^2$, $(3t-4u)^2$, $tt$; valor $(269,-83)$ hanc $(t+6u)^2+(3t+u)^2+(3t-2u)^2$; denique valor $(291,-127)$ hanc $(t+3u)^2+(3t+4u)^2+(3t-4u)^2$; singulae hae decompositiones $48$ repraesentationibus aequipollent. --- Praeter has $192$ repraesentationes autem, sive quatuor discerptiones, aliae non dabuntur, quum $770$ per nullum quadratum divisibilis sit, adeoque repraesentationes impropriae exstare non possint.

\SourceLeaf{347}
\Needspace{8\baselineskip}
\Article{290}
De formis determinantis $-1$ et $-2$, quae quibusdam exceptionibus obnoxiae erant, paucis seorsim agemus. Praemittimus observationem generalem, si $\varphi$, $\varphi'$ sint formae binariae aequivalentes quaecunque, $(\Theta)$ transformatio data illius in hanc, ex combinatione repraesentationis cuiusvis formae $\varphi$ per aliquam ternariam $f$ cum substitutione $(\Theta)$ prodire repraesentationem formae $\varphi'$ per $f$; porro ex repraesentationibus propriis ipsius $\varphi$ hoc modo oriri repraesentationes proprias formae $\varphi'$, e diversis diversas, denique e cunctis cunctas. Haec omnia per calculum facillime comprobantur. Quare una formarum $\varphi$, $\varphi'$ totidem modis per $f$ repraesentari poterit ac altera.

I. Sit primo $\varphi=tt+uu$, atque $\varphi'$ forma quaecunque alia binaria positiva det. $-1$, cui itaque $\varphi$ aequivalebit; transeat $\varphi$ in $\varphi'$ per substitutionem $t=\alpha t'+\beta u'$, $u=\gamma t'+\delta u'$. Forma $\varphi$ repraesentatur per ternariam $f=zz+yy+zz$ ponendo $x=t$, $y=u$, $z=0$; permutando $x,y,z$ hinc emergunt \textit{sex} repraesentationes, et e singulis rursus quatuor mutando signa ipsorum $t$, $u$, ita ut omnino $24$ repraesentationes diversae habeantur, quibus unica discerptio in tria quadrata aequipollet et praeter quas alias dari non posse facile perspicitur. Hinc concluditur, etiam formam $\varphi'$ unico tantum modo in tria quadrata decomponi posse, puta in $(\alpha t'+\beta u')^2$, $(\gamma t'+\delta u')^2$ et $0$, quae discerptio $24$ repraesentationibus aequivalet.

II. Sit $\varphi=tt+2uu$, $\varphi'$ quaecunque alia forma binaria positiva det. $-2$, in quam $\varphi$ transeat per substitutionem $t=\alpha t'+\beta u'$, $u=\gamma t'+\delta u'$. Tunc simili modo ut in casu praec. concluditur, $\varphi$, et proin etiam $\varphi'$, unico tantum modo in tria quadrata discerpi posse, puta $\varphi$ in $tt+uu+uu$, atque $\varphi'$ in $(\alpha t'+\beta u')^2+(\gamma t'+\delta u')^2+(\gamma t'+\delta u')^2$; talem decompositionem $24$ repraesentationibus aequipollere facile perspici potest.

Hinc colligitur, formas binarias determinantium $-1$ et $-2$ respectu multitudinis repraesentationum per ternariam $xx+yy+zz$ cum aliis formis binariis omnino convenire; quum enim in utroque casu fiat $\mu=0$, formula in art. praec. IV, tradita utique producit $24$ repraesentationes. Ratio huius rei est, quod duae exceptiones, quibus tales formae obnoxiae erant, se mutuo compensant.

Theoriam generalem repraesentationum impropriarum in art. 284 explicatam ad formam $xx+yy+zz$ applicare, brevitatis gratia supersedemus.

\SourceLeaf{348}
\Needspace{8\baselineskip}
\Article{291}
Quaestio de inveniendis omnibus repraesentationibus propriis \textit{numeri} positivi dati $M$ per formam $xx+yy+zz$ primo per art. 281 reducitur ad investigationem repraesentationum propriarum numeri $-M$ per formam $-xx-yy-zz=f$; hae vero per praecepta art. 280 ita eruuntur:

I. Evolvantur omnes classes formarum binariarum determinantis $-M$, quarum formae per $XX+YY+ZZ=F$ (cui formae ternariae adiuncta est $f$) proprie repraesentari possunt. Quando $M\equiv0,\allowbreak \ 4\text{ vel }7\ (\mathrm{mod.}\ 8)$, tales classes per art. 288 non dantur, adeoque $M$ in tria quadrata, quae divisorem communem non habeant, discerpi nequit\HistoricalNote{*)}{Haec impossibilitas etiam inde manifesta, quod summa trium quadratorum imparium necessario fit $\equiv3\ (\mathrm{mod.}\ 8)$; summa duorum imparium cum uno pari vel $\equiv2$ vel $\equiv6$; summa unius imparis cum duobus paribus vel $\equiv1$ vel $\equiv5$; denique summa trium parium vel $\equiv0$ vel $\equiv4$; sed in casu postremo repraesentatio manifesto est impropria.}. Quando vero $M\equiv1,\ 2,\ 5\text{ vel }6$, dabitur genus positivum proprie primitivum, et quando $M\equiv3$, improprie primitivum, quod omnes illas classes complectetur: designemus multitudinem harum classium per $k$.

II. Eligantur iam ex hisce classibus $k$ formae ad lubitum, e singulis una, quae sint $\varphi$, $\varphi'$, $\varphi''$ etc.; investigentur omnes omnium repraesentationes propriae per $F$. quarum itaque multitudo erit $3\, .\, 2^{\mu+3}k=K$, designante $\mu$ multitudinem factorum primorum (imparium) ipsius $M$; denique e quavis huiusmodi repraesentatione ut
\[
X=mt+nu,\quad Y=m't+n'u,\quad Z=m''t+n''u
\]
derivetur repraesentatio ipsius $M$ per $xx+yy+zz$ haec
\[
x=m'n''-m''n',\quad y=m''n-mn'',\quad z=mn'-m'n
\]
In complexu harum $K$ repraesentationum, quem per $\Omega$ designemus, omnes repraesentationes ipsius $M$ necessario contentae erunt.

III. Superest itaque tantummodo, ut inquiramus, num in $\Omega$ repraesentationes \textit{identicae} occurrere possint; et quum ex art. 280, III iam constet, eas repraesentationes in $\Omega$, quae e formis diversis \textit{e. g.} ex $\varphi$ et $\varphi'$ derivatae sint, ne%
\SourceLeaf{349}%
cessario diversas esse, sola disquisitio restat, an repraesentationes diversae eiusdem formae, \textit{e. g.} ipsius $\varphi$, per $F$, repraesentationes identicas numeri $M$ per $xx+yy+zz$ producere possint. Iam statim manifestum est, si inter repraesentationes ipsius $\varphi$ reperiatur haec
\[
X=mt+nu,\quad Y=m't+n'u,\quad Z=m''t+n''u\quad\ldots\ (r)
\]
inter easdem fore hanc
\[
X=-mt-nu,\quad Y=-m't-n'u,\quad Z=-m''t-n''u\quad\ldots\ (r')
\]
atque ex utraque derivari eandem repraesentationem ipsius $M$, quae designetur per $(R)$; examinemus itaque, num eadem $(R)$ ex aliis adhuc repraesentationibus formae $\varphi$ sequi possit. Ex art. 280, III facile deducitur, statuendo ibi $\chi=\varphi$, si omnes transformationes propriae formae $\varphi$ in se ipsam exhibeantur per
\[
t=\alpha t+\beta u,\quad u=\gamma t+\delta u
\]
omnes eas repraesentationes formae $\varphi$, e quibus $R$ sequatur, expressum iri per
\[
\begin{aligned}
x&=(\alpha m+\gamma n)t+(\beta m+\delta n)u\\
y&=(\alpha m'+\gamma n')t+(\beta m'+\delta n')u\\
z&=(\alpha m''+\gamma n'')t+(\beta m''+\delta n'')u
\end{aligned}
\]
At e theoria transformationum formarum binariarum det. negativi in art. 179 explicata sequitur, in omnibus casibus praeter $M=1$ et $M=3$, duas tantummodo transformationes proprias formae $\varphi$ in se ipsam dari, puta $\alpha,\allowbreak \beta,\allowbreak \gamma,\allowbreak \delta=1,\allowbreak 0,\allowbreak 0,\allowbreak 1$ et $=-1,0,0,-1$ resp. (quum enim $\varphi$ sit forma primitiva, id quod in art. 179 designabatur per $m$, erit vel $1$ vel $2$, et proin, praeter casus exceptos, certo $(1)$ locum ibi habebit). Quare $(R)$ e solis $r$, $r'$ provenire poterit,. adeoque quaevis repraesentatio propria numeri $M$ bis et non pluries in $\Omega$ reperietur, et multitudo omnium repraess. propriarum diversarum ipsius $M$ erit $\tfrac{1}{2}K=3\,\allowbreak  .\,\allowbreak  2^{\mu+2}k$.

Quod attinet ad casus exceptos, multitudo transformationum propriarum formae $\varphi$ in se ipsam per art. 179 erit $4$ pro $M=1$, et $6$ pro $M=3$; reveraque facile confirmatur, multitudinem repraesentationum propriarum numerorum $1$, $3$ esse $\tfrac{1}{4}K$, $\tfrac{1}{6}K$ resp.; scilicet uterque numerus unico tantum modo in tria qua%
\SourceLeaf{350}%
drata discerpi potest, $1$ in $1+0+0$, $3$ in $1+1+1$, discerptio ipsius $1$ suppeditat sex, discerptio ipsius $3$ octo repraesentationes diversas; $K$ vero fit $=24$ pro $M=1$ (ubi $\mu=0$, $k=1$) et $=48$ pro $M=3$ (ubi $\mu=1$, $k=1$).

Ceterum observamus, si $h$ designet multitudinem classium in genere principali, cui multitudo classium in quovis alio genere proprie primitivo per art. 252 aequalis est, fore $k=h$ pro $M\equiv1,\allowbreak \ 2,\allowbreak \ 5\text{ vel }6\ (\mathrm{mod.}\ 8)$, sed $k=\tfrac{1}{3}h$ pro $M\equiv3\ (\mathrm{mod.}\ 8)$, unico casu $M=3$ excepto, ubi $k=h=1$. Pro numeris itaque formae $8n+3$ multitudo repraesentationum \textit{generaliter} est $=2^{\mu+2}h$, quum in numero $3$ duae exceptiones sese compensent.

\Needspace{8\baselineskip}
\Article{292}
Discerptiones numerorum (ut formarum binariarum supra) in tria quadrata a repraesentationibus per formam $xx+yy+zz$ ita distinguimus, ut in illis ad solam quadratorum magnitudinem, in his vero insuper ad ipsorum ordinem radicumque signa respiciamus, adeoque repraesentationes $x=a$, $y=b$, $z=c$, et $x=a'$, $y=b'$, $z=c'$ pro diversis habeamus, nisi simul $a=a'$, $b=b'$, $c=c'$; discerptiones autem in $aa+bb+cc$ et in $a'a'+b'b'+c'c'$ pro una, si nullo ordinis respectu habito haec quadrata illis aequalia sunt. Hinc patet,

I. Discerptionem numeri $M$ in quadrata $aa+bb+cc$ aequipollere $48$ repraesentationibus, si nullum sit $=0$ omniaque inaequalia; $24$ autem, si \textit{vel} unum $=0$ reliqua inaequalia, \textit{vel} nullum $=0$ atque duo inter se aequalia. Si vero in discerptione numeri dati in tria quadrata duo ex his sunt $=0$, \textit{aut} unum $=0$ reliqua aequalia, \textit{aut} omnia aequalia, repraesentationibus $6$, \textit{aut} $12$, \textit{aut} $8$ aequivalens erit; sed haec evenire nequeunt nisi in casibus singularibus, ubi $M=1$ aut $2$ aut $3$ resp., siquidem repraesentationes esse debent propriae. His exclusis supponamus, multitudinem omnium discerptionum numeri $M$ in terna quadrata (divisoris communis expertia) esse $E$, atque inter has reperiri $e$ in quibus unum quadratum $0$, et $e'$ in quibus duo quadrata aequalia; illae etiam tamquam discerptiones in bina quadrata, hae tamquam discerptiones in quadratum et quadratum duplum spectari possunt. Tunc multitudo omnium repraesentationum propriarum numeri $M$ per $xx+yy+zz$ erit
\[
=24(e+e')+48(E-e-e')=48E-24(e+e')
\]
\SourceLeaf{351}
At e theoria formarum binariarum facile deducitur, $e$ fore vel $=0$ vel $=2^{\mu-1}$, prout $-1$ sit non-residuum vel residuum quadraticum ipsius $M$, nec non $e'=0$ vel $=2^{\mu-1}$, prout $-2$ non-residuum vel residuum ipsius $M$, denotante $\mu$ multitudinem factorum primorum (imparium) ipsius $M$ (v. art. 182; expositionem uberiorem hic supprimimus). Hinc facile colligitur, fore
\[
\begin{aligned}
E&=2^{\mu-2}k,&&\text{si tum $-1$ tum $-2$ sit $N.R.$ ipsius $M$;}\\
E&=2^{\mu-2}(k+2),&&\text{si uterque numerus sit residuum; denique}\\
E&=2^{\mu-2}(k+1),&&\text{si alter residuum sit alter non-residuum.}
\end{aligned}
\]
In casibus exclusis $M=1$ et $M=2$, haec formula praeberet $E=\tfrac{3}{4}$, quum esse debeat $E=1$; pro $M=3$ autem recte provenit $E=1$, exceptionibus se mutuo compensantibus.

Si itaque $M$ est numerus primus, fit $\mu=1$, adeoque $E=\tfrac{1}{2}(k+2)$ quando $M\equiv1\ (\mathrm{mod.}\ 8)$; $E=\tfrac{1}{2}(k+1)$ quando $M\equiv3$ aut $\equiv5$. Haecce theoremata specialia ab ill. Le Gendre per inductionem detecta et in commentatione egregia iam saepius laudata \textit{Hist. de l'Ac. de Paris} 1785 \textit{p.} 530 \textit{sqq.} prolata fuerunt, etsi sub forma aliquantum diversa, cuius rei ratio imprimis in eo est sita, quod aequivalentiam propriam ab impropria non distinxit, et proin classes oppositas commiscuit.

II. Ad inventionem omnium discerptionum numeri $M$ in terna quadrata (sine div. comm.) non opus est, omnes repraesentationes proprias omnium formarum $\varphi$, $\varphi'$, $\varphi''$ eruere. Primo enim facile confirmatur, omnes $(48)$ repraesentationes formae $\varphi$ ad eundem valorem expr. $\surd{}-(p,-q,r)$ pertinentes (statuendo $\varphi=(p,q,r)$) discerptionem eandem numeri $M$ praebere, adeoque sufficere, si una ex illis habeatur, sive quod eodem redit, si tantummodo omnes diversae discerptiones\HistoricalNote{*)}{Semper subintelligendum \textit{propriae}, si hanc expressionem a repraesentationibus ad discerptiones transferre lubet.} formae $\varphi$ in terna quadrata conscriptae sint, et perinde de reliquis $\varphi'$, $\varphi''$ etc. Dein si $\varphi$ est e classe non ancipite, eam formam, quae e classe opposita electa est, omnino praeterire licebit, sive e binis classibus oppositis unicam considerare sufficit. Quum enim prorsus arbitrarium sit, quaenam forma e singulis classibus eligatur, supponamus e classe opposita ei in qua est $\varphi$ eligi formam ipsi $\varphi$ oppositam, quae sit $=\varphi'$. Tunc nullo negotio perspicitur, si
